\documentclass[a4paper,12pt]{article}
\usepackage[T2A]{fontenc}
\usepackage[utf8]{inputenc}
\usepackage[russian]{babel}
\usepackage{amsmath,amssymb,mathtools}
\renewcommand{\familydefault}{\sfdefault}
\usepackage{icomma}
\usepackage[a4paper,margin=18mm]{geometry}
\usepackage{microtype}
\usepackage{tikz}
\usetikzlibrary{arrows.meta,positioning,shapes.geometric,calc}
\usepackage{tabularx,booktabs,longtable}
\usepackage{enumitem}
\usepackage{parskip}
\usepackage{fancyhdr}
\usepackage{setspace}
\usepackage{xcolor}
\usepackage{array}

\definecolor{accent}{HTML}{2A6F6B}
\definecolor{darkaccent}{HTML}{174946}
\definecolor{textgray}{HTML}{303536}
\definecolor{softgray}{HTML}{EEF2F1}
\color{textgray}
\onehalfspacing
\setlength{\parindent}{0pt}
\setlength{\headheight}{25pt}
\setlist[itemize]{leftmargin=1.5em,itemsep=0.25em,topsep=0.25em}
\setlist[enumerate]{leftmargin=1.7em,itemsep=0.55em,topsep=0.35em}
\pagestyle{fancy}
\fancyhf{}
\fancyhead[L]{\small\color{darkaccent} Практико-ориентированные задачи ОГЭ}
\fancyhead[R]{\small\color{darkaccent} 28.09.26}
\fancyfoot[C]{\small\thepage}
\renewcommand{\headrulewidth}{0.3pt}
\renewcommand{\headrule}{\hbox to\headwidth{\color{accent}\leaders\hrule height \headrulewidth\hfill}}

\newcommand{\checkitem}{\(\square\)\,}
\newcommand{\topic}[1]{\vspace{0.35em}{\large\bfseries\color{darkaccent}#1}\par\vspace{0.15em}}
\newcommand{\formula}[1]{\[\boxed{#1}\]}

\begin{document}

% ---------- TITLE ----------
\thispagestyle{empty}
\begin{center}
\vspace*{2.1cm}
{\Large\color{accent}\bfseries Чек-лист}\par
\vspace{0.55cm}
{\huge\bfseries Автомобильные шины}\par
\vspace{0.2cm}
{\Large Практико-ориентированные задачи ОГЭ}\par
\vspace{1.3cm}
{\large Ярослав Верещагин}\par
\vspace{0.35cm}
{\large 28.09.26}\par
\vfill
{\large\bfseries\color{darkaccent}Лёвин Артём Александрович}\par
\vspace{0.15cm}
{\normalsize эксклюзивно для Ярослава Верещагина}\par
\vspace{1.8cm}
\begin{tikzpicture}[remember picture,overlay]
  \draw[accent!65,line width=1.2pt]
  ($(current page.south west)+(0.7cm,1.0cm)$)
    .. controls ($(current page.south west)+(6.2cm,2.0cm)$)
    and ($(current page.south east)+(-6.0cm,0.35cm)$)
    .. ($(current page.south east)+(-0.7cm,1.25cm)$);
\end{tikzpicture}
\end{center}
\clearpage

% ---------- THEORY ----------
\section*{Главная идея}
В задачах про автомобильные шины основная сложность --- внимательно перевести маркировку и текст условия на язык величин. После этого используются обычные проценты, перевод единиц и арифметика.

\topic{1. Как читать маркировку шины}
Для маркировки вида
\[
195/65\;R15
\]
запомните смысл каждой части:
\begin{itemize}
  \item \checkitem \(B=195\) мм --- ширина шины;
  \item \checkitem \(65\) --- высота боковины в процентах от ширины, то есть \(H=0{,}65\cdot195\);
  \item \checkitem буква \(R\) указывает на радиальную конструкцию; для вычислений сама буква числового значения не даёт;
  \item \checkitem \(15\) --- диаметр диска в дюймах.
\end{itemize}

\formula{H=\frac{p}{100}\,B}
где \(p\) --- второе число маркировки.

\textbf{Пример.} Для \(225/40\;R18\):
\[
H=0{,}40\cdot225=90\text{ мм}.
\]

\topic{2. Перевод диаметра диска в миллиметры}
Один дюйм равен \(25{,}4\) мм, поэтому
\formula{d=25{,}4\,r}
где \(r\) --- число после буквы \(R\).

Например, для \(R17\):
\[
d=17\cdot25{,}4=431{,}8\text{ мм}.
\]

\clearpage
\topic{3. Полный диаметр колеса}
У колеса две боковины шины: сверху и снизу от диска. Поэтому
\formula{D=H+d+H=2H+d}

\textbf{Пример.} Для \(205/60\;R16\):
\[
H=0{,}60\cdot205=123\text{ мм},\qquad d=16\cdot25{,}4=406{,}4\text{ мм},
\]
\[
D=2\cdot123+406{,}4=652{,}4\text{ мм}.
\]

\section*{Как распознавать тип вопроса}

\begin{tabularx}{\linewidth}{>{\bfseries\raggedright\arraybackslash}p{0.28\linewidth} >{\raggedright\arraybackslash}X}
\toprule
Что спрашивают & Что делать \\
\midrule
Ширина шины & Взять первое число маркировки: \(B\). \\
Высота боковины & Найти \(H=\dfrac{p}{100}B\). \\
Диаметр диска & Найти \(d=25{,}4r\). \\
Диаметр всего колеса & Найти \(D=2H+d\). \\
Максимальная ширина при заданном диске & Найти нужный диаметр диска в таблице и сравнить только допустимые варианты. \\
Изменение диаметра в миллиметрах & Вычесть старое значение из нового: \(D_{\text{нов}}-D_{\text{стар}}\). \\
Изменение пробега за один оборот & Сравнить длины окружностей или сразу диаметры. \\
\bottomrule
\end{tabularx}

\topic{4. Пробег за один оборот колеса}
За один полный оборот колесо проходит расстояние, равное длине своей окружности:
\formula{L=\pi D}
Здесь нужен именно \textbf{полный диаметр колеса} \(D\), а не диаметр диска \(d\).

\topic{5. На сколько процентов увеличился пробег}
Если было значение \(A\), а стало \(B\), процент увеличения считается по формуле
\formula{\frac{B-A}{A}\cdot100\%}
Знаменатель --- именно \textbf{то, что было}.

Для пробега за один оборот:
\[
\frac{L_{\text{нов}}-L_{\text{стар}}}{L_{\text{стар}}}\cdot100\%
=
\frac{\pi D_{\text{нов}}-\pi D_{\text{стар}}}{\pi D_{\text{стар}}}\cdot100\%.
\]
Число \(\pi\) сокращается:
\formula{\frac{D_{\text{нов}}-D_{\text{стар}}}{D_{\text{стар}}}\cdot100\%}
Поэтому для процентного изменения пробега длину окружности отдельно вычислять не требуется.

\topic{6. Письменные вычисления и округление}
\begin{itemize}
  \item \checkitem При умножении на десятичную дробь сначала умножайте как целые числа, затем отделяйте нужное число знаков после запятой.
  \item \checkitem В процентной задаче удобно сначала умножить числитель на \(100\), а затем выполнять деление.
  \item \checkitem Не округляйте промежуточные диаметры без необходимости.
  \item \checkitem Если ответ требуется до десятых процента, итоговый процент округляйте до одного знака после запятой.
\end{itemize}

\topic{Типичные ошибки}
\begin{itemize}
  \item \checkitem Считать второе число маркировки миллиметрами: \(55\neq55\) мм, это \(55\%\) от ширины.
  \item \checkitem Забыть вторую боковину и записать \(D=H+d\) вместо \(D=2H+d\).
  \item \checkitem Смешать дюймы и миллиметры.
  \item \checkitem В формуле процентов делить на новое значение вместо исходного.
  \item \checkitem Для пробега использовать диаметр диска \(d\) вместо полного диаметра \(D\).
  \item \checkitem Округлять слишком рано и накапливать ошибку.
\end{itemize}

% ---------- FLOWCHART ----------
\clearpage
\thispagestyle{empty}
\begin{center}
{\Large\bfseries\color{darkaccent}Алгоритм: задача про маркировку шины}
\end{center}
\vspace{0.6cm}

\begin{center}
\begin{tikzpicture}[
  node distance=9mm,
  every node/.style={font=\small,align=center},
  startstop/.style={ellipse,draw=accent,line width=0.9pt,minimum width=5cm,minimum height=1.0cm},
  process/.style={rectangle,rounded corners=3pt,draw=accent,line width=0.9pt,text width=12cm,minimum height=1.15cm,inner sep=6pt},
  branch/.style={rectangle,rounded corners=3pt,draw=accent,line width=0.9pt,text width=5.8cm,minimum height=2.2cm,inner sep=6pt},
  decision/.style={diamond,aspect=2.6,draw=accent,line width=0.9pt,text width=5.7cm,inner sep=2pt},
  arrow/.style={-{Latex[length=2.6mm]},line width=0.9pt,color=darkaccent}
]
\node[startstop] (start) {Прочитайте маркировку и вопрос};
\node[process,below=of start] (decode) {Выпишите \(B\) --- ширину; \(p\) --- процент высоты; \(r\) --- диаметр диска в дюймах};
\node[process,below=of decode] (calc) {Найдите \(H=\dfrac{p}{100}B\), затем \(d=25{,}4r\); при необходимости \(D=2H+d\)};
\node[decision,below=12mm of calc] (q) {Что требуется найти?};
\node[branch,below left=13mm and 5mm of q] (direct) {Ширину, высоту, диаметр диска или полный диаметр: используйте найденную величину};
\node[branch,below right=13mm and 5mm of q] (percent) {Изменение пробега за оборот: вычислите старый и новый \(D\), затем \(\dfrac{D_{\text{нов}}-D_{\text{стар}}}{D_{\text{стар}}}\cdot100\%\)};
\node[process,below=35mm of q] (check) {Проверьте единицы, знак изменения и требуемое округление};
\node[startstop,below=of check] (finish) {Запишите только требуемый ответ};

\draw[arrow] (start) -- (decode);
\draw[arrow] (decode) -- (calc);
\draw[arrow] (calc) -- (q);
\draw[arrow] (q.west) -| node[pos=0.2,above]{обычный расчёт} (direct.north);
\draw[arrow] (q.east) -| node[pos=0.2,above]{проценты} (percent.north);
\draw[arrow] (direct.south) |- (check.west);
\draw[arrow] (percent.south) |- (check.east);
\draw[arrow] (check) -- (finish);
\end{tikzpicture}
\end{center}
\clearpage

% ---------- TRAINING ----------
\section*{Тренировочный блок --- Путь А}
\textbf{10 задач.} Решайте без калькулятора, если это соответствует условиям Вашей тренировки. В задачах на проценты итог округляйте так, как указано в условии.

\begin{enumerate}
  \item Шина имеет маркировку \(205/60\;R16\). Найдите её ширину в миллиметрах.

  \item Шина имеет маркировку \(225/40\;R18\). Найдите высоту боковины шины в миллиметрах.

  \item Диаметр диска равен \(17\) дюймам. Сколько это миллиметров?

  \item Для шины \(195/65\;R15\) найдите полный диаметр колеса в миллиметрах.

  \item На заводе установлен комплект \(215/50\;R17\). Найдите полный диаметр заводского колеса в миллиметрах.

  \item Заводские шины имеют маркировку \(185/60\;R14\), новые --- \(195/55\;R15\). На сколько миллиметров изменится полный диаметр колеса после замены?

  \item Для автомобиля разрешены следующие шины:
  \begin{center}
  \begin{tabular}{c c}
  \toprule
  Диаметр диска & Разрешённые маркировки \\
  \midrule
  15 & \(185/65\;R15\), \(195/60\;R15\), \(205/55\;R15\) \\
  16 & \(195/55\;R16\), \(205/55\;R16\), \(215/50\;R16\) \\
  17 & \(205/50\;R17\), \(215/45\;R17\), \(225/45\;R17\) \\
  \bottomrule
  \end{tabular}
  \end{center}
  Какова наибольшая ширина шины, которую можно установить при диске диаметром \(16\) дюймов?

  \item Заводские шины имеют маркировку \(185/60\;R14\), новые --- \(195/55\;R15\). На сколько процентов увеличится пробег автомобиля за один оборот колеса? Ответ округлите до десятых процента.

  \item Заводская шина имеет маркировку \(195/65\;R15\), новая --- \(205/55\;R16\). На сколько процентов изменится пробег за один оборот? Укажите, увеличится он или уменьшится. Ответ округлите до десятых процента.

  \item Заводская шина имеет маркировку \(215/50\;R17\). Рассматриваются три варианта замены: \(225/45\;R17\), \(205/55\;R17\), \(225/50\;R17\). Какой вариант даст наибольший пробег за один оборот колеса? На сколько процентов этот пробег будет больше заводского? Ответ округлите до десятых процента.
\end{enumerate}

\clearpage
% ---------- ANSWERS ----------
\section*{Ответы}
\renewcommand{\arraystretch}{1.35}
\begin{center}
\begin{tabularx}{\linewidth}{>{\centering\arraybackslash}p{1.5cm} X}
\toprule
№ & Ответ \\
\midrule
1 & \(205\) мм \\
2 & \(90\) мм \\
3 & \(431{,}8\) мм \\
4 & \(634{,}5\) мм \\
5 & \(646{,}8\) мм \\
6 & Увеличится на \(17{,}9\) мм \\
7 & \(215\) мм \\
8 & На \(3{,}1\%\) \\
9 & Уменьшится на \(0{,}4\%\) \\
10 & \(205/55\;R17\); пробег больше примерно на \(1{,}6\%\) \\
\bottomrule
\end{tabularx}
\end{center}

\vfill
\begin{center}
{\small\color{darkaccent}Лёвин Артём Александрович эксклюзивно для Ярослава Верещагина}
\end{center}

\end{document}
